% =============================================================================
% The Mathematics of Backpropagation
%
% Companion notes to llm.mojo: one section per op in the GPT-2 forward pass,
% in llm.c order. Each section states the forward computation (the anchor),
% then leaves a Derivation block to work the backward pass by hand, and a
% Result block to record the final gradient(s) once derived.
%
% Build:  make docs                      (from the repo root)
%         latexmk -pdf backprop.tex      (from docs/, equivalent)
% =============================================================================
\documentclass[journal]{IEEEtran}

\usepackage{amsmath, amssymb, amsthm}
\usepackage{mathtools}
\usepackage{bm}
\usepackage{xcolor}
\usepackage{booktabs}
\usepackage[colorlinks=true, linkcolor=blue!50!black]{hyperref}

% -----------------------------------------------------------------------------
% Macros
% -----------------------------------------------------------------------------
\newcommand{\R}{\mathbb{R}}
\newcommand{\E}{\mathbb{E}}

% Partial and total derivatives: \pdv{f}{x}, \dv{f}{x}
\newcommand{\pdv}[2]{\frac{\partial #1}{\partial #2}}
\newcommand{\dv}[2]{\frac{\mathrm{d} #1}{\mathrm{d} #2}}

% Upstream gradient of the loss w.r.t. a tensor: \dL{x} renders dL/dx.
% This is the quantity every backward kernel receives or produces.
\newcommand{\dL}[1]{\pdv{L}{#1}}

% One-hot indicator and Kronecker delta
\newcommand{\ind}[1]{\mathbb{1}\!\left[#1\right]}
\newcommand{\kron}[2]{\delta_{#1 #2}}

% softmax / layernorm / gelu as upright operators
\DeclareMathOperator{\softmax}{softmax}
\DeclareMathOperator{\layernorm}{layernorm}
\DeclareMathOperator{\gelu}{GELU}
\DeclareMathOperator{\erf}{erf}

% Workspace marker for unfinished derivations
\newcommand{\TODO}[1]{\textcolor{red!70!black}{\textbf{TODO:} #1}}

% Derivation / Result blocks
\theoremstyle{definition}
\newtheorem*{derivation}{Derivation}
\newtheorem*{result}{Result}

\title{The Mathematics of Backpropagation:\\
       Companion Derivations for \texttt{llm.mojo}}
\author{Evan~Owen}
\markboth{The Mathematics of Backpropagation}{Owen}

% =============================================================================
\begin{document}
\maketitle

\begin{abstract}
  Working derivations of the backward pass for every operation in the GPT-2
  forward pass, as implemented by the \texttt{llm.mojo} kernels. Sections are
  ordered by backward-derivation pedagogy---starting from the scalar loss and
  working toward the inputs---rather than by \texttt{llm.c} forward order, so
  that each gradient is derived only after the one it depends on. Each section
  anchors on the forward computation and leaves the derivation as structured
  workspace.
\end{abstract}

\tableofcontents

% =============================================================================
\section{Notation and Conventions}\label{sec:notation}

Shapes follow the kernel code, as collected in Table~\ref{tab:notation}.

\begin{table*}[t]
  \centering
  \caption{Symbols and their corresponding kernel parameters.}
  \label{tab:notation}
  \begin{tabular}{@{}lll@{}}
    \toprule
    Symbol & Code                         & Meaning                              \\
    \midrule
    $B$    & \texttt{batch\_size}         & batch size                           \\
    $T$    & \texttt{seq\_len}            & sequence length                      \\
    $C$    & \texttt{channels}            & model width (embedding dim)          \\
    $V$    & \texttt{vocab\_size}         & real vocabulary size                 \\
    $V_p$  & \texttt{vocab\_size\_padded} & padded vocabulary (row stride)       \\
    $H$    & \texttt{num\_heads}          & attention heads, head size $h = C/H$ \\
    \bottomrule
  \end{tabular}
\end{table*}

Indices: $b$ over batch, $t$ over positions, $i, j, k$ over channels or vocab
entries. A row of activations at one position is $x_{bt} \in \R^{C}$.

\paragraph{The backward convention.}
Every op receives the upstream gradient $\dL{y}$ of the scalar loss $L$ with
respect to its output $y$, and must produce $\dL{x}$ for each input $x$ (and
$\dL{w}$ for each parameter) via the chain rule
\begin{equation}
  \dL{x_i} \;=\; \sum_j \dL{y_j}\,\pdv{y_j}{x_i}.
  \label{eq:chain}
\end{equation}
Gradients into shared inputs \emph{accumulate} ($\mathrel{+}=$) because a
tensor used by several consumers receives a contribution from each.

\paragraph{Reading order.}
The forward pass runs encoder $\to$ layernorm $\to$ attention $\to$ matmul
$\to$ GELU $\to$ softmax $\to$ cross-entropy. The sections below instead
follow the order in which gradients are \emph{derived}: starting at the loss
(\S\ref{sec:crossentropy}) and moving back toward the embeddings
(\S\ref{sec:encoder}), each op reusing the upstream gradient established by
the section before it. \S\ref{sec:full} then reassembles these pieces into
the end-to-end backward pass.

\TODO{Decide and record: numerator vs.\ denominator layout, and where
  transposes land in the matmul gradients as a consequence.}

% =============================================================================
\section{Warm-up: Scalar Chain Rule on a Tiny Network}\label{sec:warmup}

A two-step composition to fix conventions before the tensor ops:
$z = wx + b$, $a = \sigma(z)$, $L = \tfrac{1}{2}(a - y)^2$.

\begin{derivation}
  Backward through the loss, then the sigmoid:
  \begin{align}
    \dL{a}     & = a - y                                                              \\
    \dL{z}     & = \dL{a} \cdot \pdv{a}{z}                                            \\
    \pdv{a}{z} & = \sigma(z) \cdot (1-\sigma(z))
    \intertext{and since $a = \sigma(z)$,}
    \dL{z}     & = (a - y) \cdot \pdv{a}{z} = (a - y) \cdot a \cdot (1-a).
    \intertext{For the weight:}
    \dL{w}     & = \dL{a} \cdot \pdv{a}{z} \cdot \pdv{z}{w} = \dL{z} \cdot \pdv{z}{w} \\
    \pdv{z}{w} & = x                                                                  \\
    \dL{w}     & = \dL{z} \cdot x = (a - y) \cdot a \cdot (1-a) \cdot x.
    \intertext{For the bias:}
    \dL{b}     & = \dL{a} \cdot \pdv{a}{z} \cdot \pdv{z}{b} = \dL{z} \cdot \pdv{z}{b} \\
    \pdv{z}{b} & = 1                                                                  \\
    \dL{b}     & = \dL{z} = (a - y) \cdot a \cdot (1-a).
    \intertext{For the input:}
    \dL{x}     & = \dL{a} \cdot \pdv{a}{z} \cdot \pdv{z}{x} = \dL{z} \cdot \pdv{z}{x} \\
    \pdv{z}{x} & = w                                                                  \\
    \dL{x}     & = \dL{z} \cdot w = (a - y) \cdot a \cdot (1-a) \cdot w.
  \end{align}
\end{derivation}

\begin{result}
  One chain back through the network, fanning out at the affine inputs:
  \begin{equation}
    \dL{a} \;\to\; \dL{z} \;\to\; \Bigl\{\dL{w},\; \dL{b},\; \dL{x}\Bigr\}
  \end{equation}
  with $\dL{z} = (a - y) \cdot a \cdot (1-a)$ computed once and reused:
  $\dL{w} = \dL{z} \cdot x$, \quad $\dL{b} = \dL{z}$, \quad
  $\dL{x} = \dL{z} \cdot w$.
\end{result}

% =============================================================================
\section{Cross-Entropy Loss}\label{sec:crossentropy}

Standard cross-entropy between a target distribution $y_{bt} \in \R^{V}$ and
predicted probabilities $p_{bt} \in \R^{V_p}$:
\begin{equation}
  L_{bt} \;=\; -\sum_{i<V} y_{bt,i} \log p_{bt,i}.
  \label{eq:ce-std}
\end{equation}
For language modeling the target distribution is one-hot. Writing
$\ind{P}$ for the indicator of a condition $P$,
\begin{equation}
  \ind{P} \;=\;
  \begin{cases}
    1 & P \text{ is true},  \\
    0 & P \text{ is false},
  \end{cases}
  \label{eq:indicator}
\end{equation}
the target with token $y_{bt} \in \{0,\dots,V-1\}$ is
$y_{bt,i} = \ind{i = y_{bt}}$: all probability mass on the target token,
zero elsewhere. Every term of \eqref{eq:ce-std} but one then vanishes,
and the loss collapses to the form computed by
\texttt{crossentropy\_ohe\_fwd}:
\begin{equation}
  L_{bt} \;=\; -\log p_{bt,\,y_{bt}}.
  \label{eq:ce-fwd}
\end{equation}

\subsection{Backward w.r.t.\ the probabilities}

\begin{derivation}
  Break the standard form \eqref{eq:ce-std} out term by term:
  \begin{multline}
    L_{bt} = -\bigl(
    y_{bt,0}\log p_{bt,0}
    + y_{bt,1}\log p_{bt,1} \\
    + y_{bt,2}\log p_{bt,2}
    + \cdots
    + y_{bt,V-1}\log p_{bt,V-1}
    \bigr).
    \label{eq:ce-std-expanded}
  \end{multline}
  Differentiating with respect to one coordinate $p_{bt,i}$, every term
  of the broken-out sum except the $i$-th is constant, so only its
  derivative survives:
  \begin{align}
    \dL{p_{bt,i}}
     & = -\pdv{}{p_{bt,i}}\bigl(
    y_{bt,0}\log p_{bt,0} + \cdots
    + y_{bt,i}\log p_{bt,i} \notag                   \\
     & \qquad\qquad + \cdots
    + y_{bt,V-1}\log p_{bt,V-1}
    \bigr) \notag                                    \\
     & = -\,y_{bt,i}\,\pdv{}{p_{bt,i}}\log p_{bt,i}.
    \label{eq:ce-surviving}
  \end{align}
  The surviving factor is the derivative of the logarithm,
  \begin{equation}
    \dv{}{x}\log x \;=\; \frac{1}{x}
    \qquad\Longrightarrow\qquad
    \pdv{}{p_{bt,i}}\log p_{bt,i} \;=\; \frac{1}{p_{bt,i}},
    \label{eq:dlog}
  \end{equation}
  giving the standard-sum gradient, valid for \emph{any} target
  distribution $y_{bt}$:
  \begin{equation}
    \dL{p_{bt,i}}
    \;=\;
    -\,\pdv{}{p_{bt,i}} \sum_{j<V} y_{bt,j} \log p_{bt,j}
    \;=\;
    -\frac{y_{bt,i}}{p_{bt,i}}.
    \label{eq:ce-std-bwd}
  \end{equation}
  Substituting the one-hot target $y_{bt,i} = \ind{i = y_{bt}}$,
  \eqref{eq:ce-fwd} depends only on the target coordinate $p_{bt,\,y_{bt}}$:
  \begin{equation}
    \dL{p_{bt,i}} =
    \begin{cases}
      -\dfrac{1}{p_{bt,\,y_{bt}}} & i = y_{bt},       \\
      0                           & \text{otherwise}.
    \end{cases}
  \end{equation}
\end{derivation}

\begin{result}
  \begin{equation}
    \dL{p_{bt,i}} \;=\; -\frac{\ind{i = y_{bt}}}{p_{bt,i}}.
    \label{eq:ce-bwd}
  \end{equation}
\end{result}

% =============================================================================
\section{Softmax}\label{sec:softmax}

Forward: from logits $\ell_{bt} \in \R^{V_p}$ (real entries $i < V$),
\begin{equation}
  p_{bt,i} \;=\; \frac{e^{\ell_{bt,i} - m_{bt}}}{\sum_{j<V} e^{\ell_{bt,j} - m_{bt}}},
  \qquad m_{bt} = \max_{j<V} \ell_{bt,j}.
  \label{eq:softmax-fwd}
\end{equation}

\subsection{The Jacobian}

Row indices $bt$ are dropped throughout: each row is an independent
softmax, so the Jacobian is per-row.

\begin{derivation}
  The object here is the Jacobian entry $\pdv{p_i}{\ell_j}$ alone: how
  one probability moves with one logit. The loss enters only in the
  subsections that follow, where an upstream gradient is chained through
  these entries. Before differentiating: the max-subtraction in
  \eqref{eq:softmax-fwd} cancels, because for any constant $c$
  \begin{equation}
    \frac{e^{\ell_i - c}}{\sum_{j<V} e^{\ell_j - c}}
    \;=\;
    \frac{e^{-c}\,e^{\ell_i}}{e^{-c}\sum_{j<V} e^{\ell_j}}
    \;=\;
    \frac{e^{\ell_i}}{\sum_{j<V} e^{\ell_j}},
    \label{eq:softmax-shift}
  \end{equation}
  so $m_{bt}$ does not change the derivative and the unshifted form can be
  differentiated. Break the denominator sum out in the same format:
  \begin{equation}
    p_i = \frac{e^{\ell_i}}{S},
    \qquad
    S = e^{\ell_0} + e^{\ell_1} + e^{\ell_2} + \cdots + e^{\ell_{V-1}}.
    \label{eq:softmax-S}
  \end{equation}
  Every entry comes from the quotient rule: for $u/v$ and a scalar
  argument $x$,
  \begin{equation}
    \pdv{}{x}\,\frac{u}{v}
    \;=\;
    \frac{\pdv{u}{x}\,v \;-\; u\,\pdv{v}{x}}{v^{2}}.
    \label{eq:quotient}
  \end{equation}
  Work the $2\times2$ block of entries in $p_1, p_2$ and
  $\ell_1, \ell_2$ explicitly before generalizing.

  \paragraph{Entry $\pdv{p_1}{\ell_1}$}
  Here $u = e^{\ell_1}$ and $v = S$, and \emph{both} move with $\ell_1$:
  \begin{align}
    \pdv{u}{\ell_1} & = \pdv{}{\ell_1}\,e^{\ell_1} = e^{\ell_1},
    \label{eq:q11-u}                                                \\
    \pdv{v}{\ell_1} & = \pdv{}{\ell_1}\bigl(e^{\ell_0} + e^{\ell_1}
    + e^{\ell_2} + \cdots + e^{\ell_{V-1}}\bigr) = e^{\ell_1},
    \label{eq:q11-v}
  \end{align}
  every term of the broken-out sum except $e^{\ell_1}$ being constant.
  Substituting \eqref{eq:q11-u} and \eqref{eq:q11-v} into the quotient
  rule \eqref{eq:quotient} with the values filled in:
  \begin{align}
    \pdv{p_1}{\ell_1}
     & = \frac{e^{\ell_1}\bigl(e^{\ell_0} + e^{\ell_1} + \cdots
      + e^{\ell_{V-1}}\bigr)
      - e^{\ell_1}\,e^{\ell_1}}
    {\bigl(e^{\ell_0} + e^{\ell_1} + \cdots + e^{\ell_{V-1}}\bigr)^{2}}
    \notag                                                      \\
     & = \frac{e^{\ell_1}\,S - e^{\ell_1}\,e^{\ell_1}}{S^{2}}
    = \frac{e^{\ell_1}}{S}
    - \left(\frac{e^{\ell_1}}{S}\right)^{\!2} \notag            \\
     & = p_1 - p_1^{2} \;=\; p_1\,(1 - p_1).
    \label{eq:q11}
  \end{align}

  \paragraph{Entry $\pdv{p_1}{\ell_2}$}
  Now the numerator does \emph{not} move: $u = e^{\ell_1}$ is constant
  in $\ell_2$, so $\pdv{u}{\ell_2} = 0$, while again exactly one term of
  the broken-out denominator survives:
  \begin{equation}
    \pdv{v}{\ell_2} = \pdv{}{\ell_2}\bigl(e^{\ell_0} + e^{\ell_1}
    + e^{\ell_2} + \cdots + e^{\ell_{V-1}}\bigr) = e^{\ell_2}.
    \label{eq:q12-v}
  \end{equation}
  Substituting into \eqref{eq:quotient} with the values filled in:
  \begin{align}
    \pdv{p_1}{\ell_2}
     & = \frac{0 \cdot \bigl(e^{\ell_0} + e^{\ell_1} + \cdots
      + e^{\ell_{V-1}}\bigr)
      - e^{\ell_1}\,e^{\ell_2}}
    {\bigl(e^{\ell_0} + e^{\ell_1} + \cdots + e^{\ell_{V-1}}\bigr)^{2}}
    \notag                                                    \\
     & = -\,\frac{e^{\ell_1}}{S}\cdot\frac{e^{\ell_2}}{S}
    \;=\; -\,p_1\,p_2.
    \label{eq:q12}
  \end{align}

  \paragraph{Entries $\pdv{p_2}{\ell_1}$ and $\pdv{p_2}{\ell_2}$}
  The same two computations with the indices swapped. For
  $\pdv{p_2}{\ell_1}$: $u = e^{\ell_2}$ is constant in $\ell_1$ and
  $\pdv{v}{\ell_1} = e^{\ell_1}$ by \eqref{eq:q11-v}; for
  $\pdv{p_2}{\ell_2}$ both numerator and denominator move, exactly as in
  \eqref{eq:q11}:
  \begin{align}
    \pdv{p_2}{\ell_1}
     & = \frac{0 \cdot S - e^{\ell_2}\,e^{\ell_1}}{S^{2}}
    = -\,p_2\,p_1,
    \label{eq:q21}                                            \\
    \pdv{p_2}{\ell_2}
     & = \frac{e^{\ell_2}\,S - e^{\ell_2}\,e^{\ell_2}}{S^{2}}
    = p_2\,(1 - p_2).
    \label{eq:q22}
  \end{align}

  \paragraph{The pattern}
  Arranged as a block of the Jacobian, the four entries are
  \begin{equation}
    \begin{pmatrix}
      \pdv{p_1}{\ell_1} & \pdv{p_1}{\ell_2} \\[6pt]
      \pdv{p_2}{\ell_1} & \pdv{p_2}{\ell_2}
    \end{pmatrix}
    =
    \begin{pmatrix}
      p_1\,(1-p_1) & -\,p_1\,p_2  \\[2pt]
      -\,p_2\,p_1  & p_2\,(1-p_2)
    \end{pmatrix}:
    \label{eq:softmax-block}
  \end{equation}
  diagonal entries $p_i\,(1 - p_i)$, off-diagonal entries
  $-\,p_i\,p_j$. The general entry repeats the same quotient rule. Each
  $v$-derivative above is one instance of the general fact that exactly
  one term of the broken-out $S$ involves $\ell_j$:
  \begin{equation}
    \pdv{S}{\ell_j}
    = \pdv{}{\ell_j}\bigl(e^{\ell_0} + \cdots + e^{\ell_j} + \cdots
    + e^{\ell_{V-1}}\bigr)
    = e^{\ell_j}.
    \label{eq:softmax-dS}
  \end{equation}
  The diagonal case $j = i$ repeats \eqref{eq:q11} (the numerator also
  moves):
  \begin{align}
    \pdv{p_i}{\ell_i}
     & = \frac{e^{\ell_i}\,S - e^{\ell_i}\,e^{\ell_i}}{S^2}
    = \frac{e^{\ell_i}}{S}
    - \left(\frac{e^{\ell_i}}{S}\right)^{\!2} \notag        \\
     & = p_i - p_i^2 = p_i\,(1 - p_i).
    \label{eq:softmax-diag}
  \end{align}
  The off-diagonal case $j \neq i$ repeats \eqref{eq:q12} (the numerator
  $e^{\ell_i}$ is constant in $\ell_j$):
  \begin{equation}
    \pdv{p_i}{\ell_j}
    = \frac{0 \cdot S - e^{\ell_i}\,e^{\ell_j}}{S^2}
    = -\,p_i\,p_j.
    \label{eq:softmax-offdiag}
  \end{equation}
\end{derivation}

\begin{result}
  The two cases combine into one expression via the Kronecker delta, the
  standard notation for an entry of the identity matrix. It is the indicator
  \eqref{eq:indicator} specialized to an index equality,
  $\kron{i}{j} = \ind{i = j}$:
  \begin{equation}
    \kron{i}{j} \;=\;
    \begin{cases}
      1 & i = j,    \\
      0 & i \neq j,
    \end{cases}
    \label{eq:kron}
  \end{equation}
  which selects the diagonal case:
  \begin{equation}
    \pdv{p_i}{\ell_j} \;=\; p_i\,\bigl(\kron{i}{j} - p_j\bigr).
    \label{eq:softmax-jac}
  \end{equation}
  Check against \eqref{eq:softmax-diag} and \eqref{eq:softmax-offdiag}:
  $j = i$ gives $p_i\,(1 - p_i)$, and $j \neq i$ gives $-\,p_i\,p_j$.
\end{result}

\subsection{Backward w.r.t.\ the logits}

\begin{derivation}
  The upstream gradient $\dL{p}$ enters through the chain rule
  \eqref{eq:chain} summed over the row, one full column of the
  Jacobian \eqref{eq:softmax-jac} per output $i$. Broken out:
  \begin{align}
    \dL{\ell_i}
     & = \sum_{j<V} \dL{p_j}\,\pdv{p_j}{\ell_i} \notag      \\
     & = \dL{p_0}\pdv{p_0}{\ell_i}
    + \dL{p_1}\pdv{p_1}{\ell_i}
    + \dL{p_2}\pdv{p_2}{\ell_i} \notag                      \\
     & \qquad + \cdots + \dL{p_{V-1}}\pdv{p_{V-1}}{\ell_i}.
    \label{eq:softmax-chain}
  \end{align}
  Substitute \eqref{eq:softmax-diag} and \eqref{eq:softmax-offdiag}; the
  $j = i$ term is the diagonal case, every other term is off-diagonal:
  \begin{align}
    \dL{\ell_i}
     & = \dL{p_i}\,p_i\,(1 - p_i)
    + \sum_{j \neq i} \dL{p_j}\,(-\,p_j\,p_i) \notag   \\
     & = p_i\,\dL{p_i}
    - p_i\,p_i\,\dL{p_i}
    - p_i \sum_{j \neq i} \dL{p_j}\,p_j.
    \intertext{The middle term is exactly the missing $j = i$ element of
      the sum, so it re-absorbs:}
    \dL{\ell_i}
     & = p_i\,\dL{p_i} - p_i \sum_{j<V} \dL{p_j}\,p_j,
  \end{align}
  where the remaining sum is the same one for every $i$ in the row.
  Broken out in the same format, it is
  \begin{multline}
    \sum_{j<V} \dL{p_j}\,p_j
    = \dL{p_0}\,p_0 + \dL{p_1}\,p_1 \\
    + \dL{p_2}\,p_2 + \cdots + \dL{p_{V-1}}\,p_{V-1}.
    \label{eq:softmax-dot}
  \end{multline}
\end{derivation}

\begin{result}
  \begin{equation}
    \dL{\ell_i} \;=\;
    p_i \left(\dL{p_i} - \sum_{j<V} \dL{p_j}\,p_j\right),
    \qquad i < V,
    \label{eq:softmax-bwd}
  \end{equation}
  and $\dL{\ell_i} = 0$ for the padded region $i \geq V$. Summing
  \eqref{eq:softmax-bwd} over the row, the leading $p_i$ distributes into
  both terms; the probabilities multiplying the (constant) dot product
  sum to one, leaving two copies of the same dot product:
  \begin{align}
    \sum_{i<V} \dL{\ell_i}
     & = \sum_{i<V} p_i\,\dL{p_i}
    - \Bigl(\sum_{i<V} p_i\Bigr)
    \sum_{j<V} \dL{p_j}\,p_j \notag \\
     & = \sum_{i<V} p_i\,\dL{p_i}
    - \sum_{j<V} \dL{p_j}\,p_j
    \;=\; 0.
    \label{eq:softmax-zerosum}
  \end{align}
  The gradient row always sums to zero, consistent with the shift
  invariance \eqref{eq:softmax-shift}.
\end{result}

\subsection{Fused softmax + cross-entropy}

The composition is the case that matters for the model: backprop through
\eqref{eq:ce-fwd} and \eqref{eq:softmax-fwd} together and watch the
Jacobian collapse.

\begin{derivation}
  Write out the chain per logit coordinate, full indices restored
  (upstream index $i$ over the probabilities, $j$ over the logits):
  \begin{equation}
    \dL{\ell_{bt,j}}
    \;=\;
    \sum_{i<V} \dL{p_{bt,i}}\,\pdv{p_{bt,i}}{\ell_{bt,j}}.
    \label{eq:fused-chain}
  \end{equation}
  Substitute what is already known into the combined equation: the
  cross-entropy gradient \eqref{eq:ce-bwd} for the first factor and the
  Jacobian \eqref{eq:softmax-jac} for the second:
  \begin{align}
    \dL{\ell_{bt,j}}
     & = \sum_{i<V}
    \left(-\frac{\ind{i = y_{bt}}}{p_{bt,i}}\right)
    p_{bt,i}\,\bigl(\kron{i}{j} - p_{bt,j}\bigr)
    \intertext{The $p_{bt,i}$ in the denominator cancels against the
      $p_{bt,i}$ the Jacobian contributes:}
     & = -\sum_{i<V} \ind{i = y_{bt}}\,
    \bigl(\kron{i}{j} - p_{bt,j}\bigr)
    \intertext{and the one-hot indicator keeps exactly one term of the
      sum, $i = y_{bt}$:}
     & = -\bigl(\kron{y_{bt}}{j} - p_{bt,j}\bigr)
    = p_{bt,j} - \ind{j = y_{bt}}.
    \label{eq:fused-collapse}
  \end{align}
\end{derivation}

\begin{result}
  \begin{equation}
    \dL{\ell_{bt,j}} \;=\; p_{bt,j} - \ind{j = y_{bt}},
    \qquad j < V,
    \label{eq:fused-bwd}
  \end{equation}
  and $\dL{\ell_{bt,j}} = 0$ for the padded region $j \geq V$. The
  Jacobian collapsed entirely: no sum over the row survives in
  \eqref{eq:fused-bwd}.

  Cross-check via the general backward \eqref{eq:softmax-bwd}: with the
  cross-entropy upstream \eqref{eq:ce-bwd}, the dot product evaluates to
  \begin{equation}
    \sum_{j<V} \dL{p_{bt,j}}\,p_{bt,j}
    = \sum_{j<V} \left(-\frac{\ind{j = y_{bt}}}{p_{bt,j}}\right) p_{bt,j}
    = -1,
    \label{eq:fused-dot}
  \end{equation}
  a constant: for a one-hot target the row-wide weighted sum is known
  before reading a single element. Substituting $-1$ for the dot product
  in \eqref{eq:softmax-bwd}, the $p_{bt,i}$ outside the parentheses
  cancels the reciprocal inside them:
  \begin{align}
    \dL{\ell_{bt,i}}
     & = p_{bt,i}\left(\dL{p_{bt,i}} - (-1)\right) \notag      \\
     & = p_{bt,i}
    \left(-\frac{\ind{i = y_{bt}}}{p_{bt,i}} + 1\right) \notag \\
     & = p_{bt,i} - \ind{i = y_{bt}},
    \label{eq:fused-check}
  \end{align}
  agreeing with \eqref{eq:fused-bwd}: chaining through the Jacobian
  directly and substituting into the general backward are two routes to
  the same gradient. The zero-sum check of \S\ref{sec:softmax} also
  holds: the probabilities sum to one, and the indicator fires exactly
  once per row:
  \begin{equation}
    \sum_{j<V} \bigl(p_{bt,j} - \ind{j = y_{bt}}\bigr)
    \;=\; 1 - 1 \;=\; 0.
    \label{eq:fused-zerosum}
  \end{equation}

  The collapse did not need one-hotness. It needed only the standard-sum
  gradient \eqref{eq:ce-std-bwd} and the fact that the target is a
  distribution, $\sum_{i<V} y_{bt,i} = 1$. Repeating
  \eqref{eq:fused-chain} with a general target distribution,
  \begin{align}
    \dL{\ell_{bt,j}}
     & = \sum_{i<V}
    \left(-\frac{y_{bt,i}}{p_{bt,i}}\right)
    p_{bt,i}\,\bigl(\kron{i}{j} - p_{bt,j}\bigr) \notag \\
     & = -\,y_{bt,j} + p_{bt,j}\sum_{i<V} y_{bt,i}
    \;=\; p_{bt,j} - y_{bt,j},
    \label{eq:fused-general}
  \end{align}
  the one-hot result \eqref{eq:fused-bwd} being the special case
  $y_{bt,j} = \ind{j = y_{bt}}$. Softmax + cross-entropy always
  backpropagates as \emph{predicted minus target}.
\end{result}

% =============================================================================
\section{Linear Layer (Matmul)}\label{sec:matmul}

Forward: $Y = XW^{\top} + b$ with $X \in \R^{(B\cdot T)\times C_{\text{in}}}$,
$W \in \R^{C_{\text{out}}\times C_{\text{in}}}$, $b \in \R^{C_{\text{out}}}$.

\begin{derivation}
  \TODO{Derive $\dL{X}$, $\dL{W}$, $\dL{b}$ from \eqref{eq:chain}; check
    every shape. Note where the reduction over $B\cdot T$ appears.}
\end{derivation}

\begin{result}
  \TODO{$\dL{X} = \;?$, \quad $\dL{W} = \;?$, \quad $\dL{b} = \;?$}
\end{result}

% =============================================================================
\section{GELU}\label{sec:gelu}

Forward (tanh approximation, as in llm.c):
\begin{equation}
  \gelu(x) \;=\; \tfrac{1}{2}\,x\left(1 + \tanh\!\left[\sqrt{2/\pi}\,
      (x + 0.044715\,x^{3})\right]\right).
  \label{eq:gelu-fwd}
\end{equation}

\begin{derivation}
  \TODO{Differentiate \eqref{eq:gelu-fwd}. Elementwise, so the chain rule is
    a pointwise product.}
\end{derivation}

\begin{result}
  \TODO{$\gelu'(x) = \;?$}
\end{result}

% =============================================================================
\section{LayerNorm}\label{sec:layernorm}

Forward, per position $x \in \R^{C}$ with scale $\gamma$ and shift $\beta$:
\begin{equation}
  \begin{aligned}
    \mu & = \tfrac{1}{C}\sum_i x_i,
    \qquad
    \sigma^2 = \tfrac{1}{C}\sum_i (x_i - \mu)^2,                                 \\
    y_i & = \gamma_i\,\frac{x_i - \mu}{\sqrt{\sigma^2 + \varepsilon}} + \beta_i.
  \end{aligned}
  \label{eq:ln-fwd}
\end{equation}

\begin{derivation}
  \TODO{The interesting one: $\mu$ and $\sigma^2$ both depend on every
    $x_j$, so $\pdv{y_i}{x_j}$ has three terms. Work it via
    $\hat{x}_i = (x_i - \mu)/\sqrt{\sigma^2 + \varepsilon}$.}
\end{derivation}

\begin{result}
  \TODO{$\dL{x_j} = \;?$, \quad $\dL{\gamma_i} = \;?$, \quad $\dL{\beta_i} = \;?$}
\end{result}

% =============================================================================
\section{Attention}\label{sec:attention}

Forward, per head: $A = \softmax\!\big(QK^{\top}/\sqrt{h} + M\big)$,
$Y = AV$, with causal mask $M$ and $Q = XW_Q$, etc.

\begin{derivation}
  \TODO{Backprop in stages: $\dL{V}$ and $\dL{A}$ from $Y = AV$; then
    $\dL{(QK^{\top})}$ through the row-wise softmax (reuse
    \S\ref{sec:softmax}); then $\dL{Q}$, $\dL{K}$. Track $\sqrt{h}$ and the
    mask.}
\end{derivation}

\begin{result}
  \TODO{Gradient for each of $Q$, $K$, $V$.}
\end{result}

% =============================================================================
\section{Embeddings (Encoder)}\label{sec:encoder}

Forward: $x_{bt} = W_{E}[\,\mathrm{tok}_{bt}\,] + W_{P}[\,t\,]$.

\begin{derivation}
  \TODO{Gradients of a gather are a scatter-add; multiple $(b,t)$ hitting
    the same token row accumulate. Connect to the $\mathrel{+}=$ convention
    of \S\ref{sec:notation}.}
\end{derivation}

\begin{result}
  \TODO{$\dL{W_E} = \;?$, \quad $\dL{W_P} = \;?$}
\end{result}

% =============================================================================
\section{The Full Backward Pass}\label{sec:full}

Compose \S\ref{sec:crossentropy}--\S\ref{sec:encoder} in reverse forward
order, including the residual stream (where gradients add) and weight tying
between $W_E$ and the LM head.

\TODO{Write the end-to-end gradient flow for one transformer block, then
  the whole model.}

% =============================================================================
\appendices
\section{Gradient Checking}\label{app:gradcheck}

Numerical check used to validate every derivation above:
\begin{equation}
  \pdv{L}{\theta_i} \;\approx\;
  \frac{L(\theta + \epsilon e_i) - L(\theta - \epsilon e_i)}{2\epsilon}.
\end{equation}

\TODO{Record per-dtype tolerances and how they relate to
  \texttt{tests/\_dtypes.py}.}

\section{Useful Identities}\label{app:identities}

\begin{itemize}
  \item $\dv{}{x}\log x = 1/x$; \quad $\dv{}{x} e^{x} = e^{x}$; \quad
        $\dv{}{x}\tanh x = 1 - \tanh^2 x$.
  \item log-sum-exp and its gradient: $\nabla \log \sum_j e^{x_j} = \softmax(x)$.
  \item \TODO{Collect identities as they get used.}
\end{itemize}

\end{document}
